% Generated draft: requires mathematical and source review.
\documentclass[11pt]{article}
\usepackage[margin=25mm]{geometry}
\usepackage{fontspec}
\setmainfont{Latin Modern Roman}
\usepackage{amsmath,amssymb,mathtools,unicode-math}
\setmathfont{Latin Modern Math}
\usepackage{enumitem,needspace,xurl,hyperref}
\setlength{\parindent}{0pt}
\setlength{\parskip}{5pt}
\setlength{\emergencystretch}{3em}
\newcommand{\N}{\mathbb{N}}
\newcommand{\Z}{\mathbb{Z}}
\newcommand{\Q}{\mathbb{Q}}
\newcommand{\R}{\mathbb{R}}
\newcommand{\C}{\mathbb{C}}
\newcommand{\F}{\mathbb{F}}
\newcommand{\A}{\mathbb{A}}
\newcommand{\PP}{\mathbb{P}}
\newcommand{\E}{\mathbb{E}}
\newcommand{\eps}{\varepsilon}
\newcommand{\dd}{\,\mathrm{d}}
\begin{document}
\section*{UnsolvedMath ID: 30001036}
\textbf{Source problem code:} catalogue entry
\subsection{Definitions and mathematical statement}
For a finitely generated group $G$ with a finite symmetric generating set, let $d$ be its Cayley-graph distance from the identity. A standard Floyd construction uses a positive summable function $f:\N\to(0,\infty)$ and gives an edge $\{u,v\}$ length $f(\min\{d(1,u),d(1,v)\})$. The induced path metric on vertices has a completion $\overline G_f$, and the Floyd boundary is $\partial_fG=\overline G_f\setminus G$. Standard regularity conditions often include monotonicity and a uniform lower bound on $f(n+1)/f(n)$; the source instead refers to its own unstated standard conditions, which must be recovered.

Retain all three questions from the source:
\begin{enumerate}
\item For the intended allowed Floyd function, does there exist $G$ with nontrivial $\partial_fG$ that is not relatively hyperbolic with respect to any allowed family of finitely generated peripheral subgroups?
\item For the particular infinitely presented group $H$ constructed in the report, is $\partial_fH$ trivial?
\item Is that same $H$ not relatively hyperbolic with respect to every such family?
\end{enumerate}
The second and third are separate questions about a specified but unrecovered group; no answers are inferred from its description as a candidate.

For a substantive relative-hyperbolicity question, peripheral structures must exclude the vacuous choice $\{G\}$; the usual setting uses proper peripheral subgroups. The extract does not specify whether the family is finite or which degeneracies are permitted. Likewise ``nontrivial Floyd boundary'' may mean at least three points in the conventional nonelementary sense, rather than merely nonempty; the exact cardinality convention is absent.

The original keeps $f$ subject to the report's conditions, while the clean statement suppresses it as a hypothesis. Its quantifier (fixed, arbitrary allowed, or existential) and the finite-generation/generating-set conventions must therefore be confirmed. The group $H$ has no presentation or construction in this record and is not replaced by another group mentioned in the workshop.
\subsection{Short English statement}
Seek a nontrivial Floyd boundary without any permissible relatively hyperbolic structure, and decide two separate boundary/hyperbolicity questions for the report's infinitely presented group. The Floyd hypotheses, peripheral conventions and candidate group are not supplied.
\subsection{Sources}
\begin{itemize}
\item Supplied report: \emph{Geometric Group Theory, Hyperbolic Dynamics and Symplectic Geometry} (2009), Oberwolfach Reports, DOI 10.4171/owr/2008/42, pp. 2337--2406; proposed year 2008. \url{https://doi.org/10.4171/owr/2008/42}. The verification identifies Questions 6--8.
\item Supplied Yang reference on peripheral structures: \url{https://arxiv.org/abs/1101.2990}.
\item Supplied Mahan reference on controlled Floyd separation: \url{https://mathweb.tifr.res.in/~mahan/floyd.pdf}.
\item Local UnsolvedMath registry, permanent ID 30001036, catalogue entry; created 25 August 2026, updated 16 September 2026. Its corrected\_verified note says the general existence and both candidate questions are preserved. Literature checked at \texttt{2026-08-21}: the assessment is supplied dated curation, not independently verified. It labels the general nontrivial-Floyd-implies-relative-hyperbolicity conjecture unresolved while mentioning counterexamples to weaker formulations; their precise versions are not provided.
\end{itemize}
\end{document}
